\documentclass[12pt]{article}

\usepackage[margin=2cm]{geometry}
\usepackage{graphicx}
\usepackage{caption}
\usepackage[numbers]{natbib}
\usepackage{minted}   % syntax highlighting in code blocks
\usepackage{listings}
\usepackage{xcolor}
\usepackage[colorlinks=true, linkcolor=blue, urlcolor=blue, citecolor=blue]{hyperref}
\usepackage{microtype}  % nicer typesetting
\usepackage{libertinus}
\usepackage{inconsolata}  % nicer font for comments in code blocks
\usepackage{texshade}
\usepackage{tikz}
\usetikzlibrary{shapes, arrows, positioning} 
\usepackage{smartdiagram}
\usesmartdiagramlibrary{additions}

\linespread{1.05}

\title{\LARGE\textbf{Short Analysis Pipeline}\\\vspace{1mm}\large BIOLM0051}
\author{Archie Benn}
\date{}

% COMPILE WITH THE FOLLOWING: 
%pdflatex -shell-escape report.tex 
\begin{document}
\maketitle

%%%%%%%%%%%%%%%%%%%%%%%%%%%%%%%%
\subsection*{Scenario}
A suitcase containing several packages of unidentified frozen meat has been intercepted at Heathrow Airport by UK Border Force officers during routine inspections. The shipment had no import documentation, raising concerns that the meat may originate from protected or endangered species.\newline

Given the serious implications for public health, food safety, and wildlife conservation, the case has been referred to the Animal and Plant Health Agency (APHA). As a bioinformatician assisting with the inquiry, you have been tasked with determining the species of origin for each meat sample. DNA barcoding has already been performed in the laboratory, and you have now received the resulting FASTAQ sequence files for analysis.
%%%%%%%%%%%%%%%%%%%%%%%%%%%%%%%%



%%%%%%%%%%%%%%%%%%%%%%%%%%%%%%%%
\section*{Abstract} %200 WORDS

%%%%%%%%%%%%%%%%%%%%%%%%%%%%%%%%



%%%%%%%%%%%%%%%%%%%%%%%%%%%%%%%%
\section*{Introduction}

%%%%%%%%%%%%%%%%%%%%%%%%%%%%%%%%



%%%%%%%%%%%%%%%%%%%%%%%%%%%%%%%%
\section*{Methods}
The FASTQ files in this report are small and the use of HPC resources in this case is not necessary. However, this workflow is designed to scale easily to similar analyses with much larger datasets. All steps here are able to run on an HPC cluster without modification, demonstrating reproduciility, automation, and capacity for high-throughput analyses.

%%%%%%%%%%%%%%%
% Pipeline workflow figure
%%%%%%%%%%%%%%%
\begin{figure}[htbp]
\begin{center}
\vspace{3em}
\smartdiagramset{
    set color list={red!20, red!30,red!40,red!50},
    sequence item border size=1.2\pgflinewidth,
    sequence item border color=black,sequence item font size=\footnotesize,
    sequence item text color=black, 
    additions={
            additional item shape=rectangle,
            additional item fill color=black!10,
            additional item border color=black,
            additional item font=\footnotesize,
            additional arrow line width=2pt,
            additional arrow tip=to,
            additional arrow color=black,
        }
}
\smartdiagramadd[sequence diagram]{Concatenate FASTQ,ORF finder,BLAST, Alignment}
{above of sequence-item1/FASTQ files, below of sequence-item4/Aligned Sequences}
\smartdiagramconnect{to-}{sequence-item1/additional-module1}
\smartdiagramconnect{-to}{sequence-item4/additional-module2}
\vspace{4em}
\caption{Workflow of Snakemake analysis pipeline.}
    \label{fig:workflow}
\end{center}
\end{figure}
%%%%%%%%%%%%%%%
% Pipeline workflow figure
%%%%%%%%%%%%%%%




Some python code: 
%%%%%%%%%%%%%%%
% Python block
%%%%%%%%%%%%%%%
\begin{minted}[fontsize=\small, bgcolor=gray!1, framesep=1.5mm, frame=single]{python}
from Bio import SeqIO
for record in SeqIO.parse("sequence.fasta", "fasta"):
    print(record.id, record.seq.translate())
\end{minted} 
%%%%%%%%%%%%%%%
% Python block
%%%%%%%%%%%%%%%

Some Bash script:
%%%%%%%%%%%%%%%
% Bash block
%%%%%%%%%%%%%%%
\begin{minted}[fontsize=\small, frame=single, framesep=1.5mm, bgcolor=gray!1]{bash}
#!/bin/bash
#SBATCH --job-name=blast_job
#SBATCH --partition=teach_cpu
#SBATCH --time=01:00:00
#SBATCH --mem=2000M

blastn -db nt -query sequence.fa -out results.out -remote
\end{minted}
%%%%%%%%%%%%%%%
% Bash block
%%%%%%%%%%%%%%%
%%%%%%%%%%%%%%%%%%%%%%%%%%%%%%%%



%%%%%%%%%%%%%%%%%%%%%%%%%%%%%%%%
\section*{Results}
In this study, we analysed the gene expression levels across multiple samples using RNA-seq data. The expression values $x_i$ for gene $i$ were normalised using the log-transformation:
\[
y_i = \log_2(x_i + 1)
\]
to stabilise variance. Data processing was performed using Python scripts:

\begin{minted}[frame=single, framesep=1.5mm, bgcolor=gray!1]{python}
import pandas as pd
data = pd.read_csv("expression.csv")
normalized = data.apply(lambda x: np.log2(x+1))
\end{minted}

Alignment results using Clustal Omega on some test data: 
%%%%%%%%%%%%%%%
% Texshade alignment
%%%%%%%%%%%%%%%
\begin{figure}[htbp]
\begin{texshade}{../data/testing/clustalo-I20251025-093857-0026-94127824-p1m.aln-msf}
\threshold[80]{50}
\setends{1}{0..110}
\showconsensus[ColdHot]{bottom}
\defconsensus{.}{lower}{upper}
\showlegend
\end{texshade}
\caption{Alignment using Clustal Omega and Identity mode in texshade}
\label{fig:workflow}
\end{figure}
%%%%%%%%%%%%%%%
% Texshade alignment
%%%%%%%%%%%%%%%
%%%%%%%%%%%%%%%%%%%%%%%%%%%%%%%%



%%%%%%%%%%%%%%%%%%%%%%%%%%%%%%%%
\section*{Discussion}
Yadadadada \citep{kent2002blat} and then also \citep{wolin1997microbe} and maybe even from this guy \cite{baker1997signaling,wolin1997microbe}
%%%%%%%%%%%%%%%%%%%%%%%%%%%%%%%%



%%%%%%%%%%%%%%%%%%%%%%%%%%%%%%%% check .bib ordering for years and surnames etc.
\bibliographystyle{unsrtnat}
\bibliography{references}
%%%%%%%%%%%%%%%%%%%%%%%%%%%%%%%%




%Acknowledging the use of HPC Facilities
%When you write articles about your research for publication, 
%conference proceedings or for other reasons, please acknowledge use of the ACRC facilities by including the following:
%%%%%%%%%%%%%%%%%%%%%%%%%%%%%%%%
\section*{Acknowledgments}
This work was carried out using the computational facilities of the Advanced Computing Research Centre, University of Bristol - \url{http://www.bristol.ac.uk/acrc/}
%%%%%%%%%%%%%%%%%%%%%%%%%%%%%%%%




%%%%%%%%%%%%%%%%%%%%%%%%%%%%%%%%
\section*{Supplementary Data}
Supplementary data can be found within the GitHub repository for this project: \url{https://github.com/archiebenn/BIOLM0051_analysis.git}
%%%%%%%%%%%%%%%%%%%%%%%%%%%%%%%%
\end{document}